% !TEX root = ../main.tex
% §7 Results. Main table: results/stage6_main/summary.csv and latency.csv; significance:
% results/stage8_analysis/significance.csv (src/analyze.py, N_BOOT = 1000); ablations:
% results/stage7_ablation/ablation_summary.csv.
\section{Results}
\label{sec:results}

\begin{table*}[t]
  \centering
  \footnotesize
  \setlength{\tabcolsep}{3pt}
  \begin{tabular}{@{}lcccccccc@{}}
    \toprule
    Method & PR-AUC & ROC-AUC & F1 & F2 & FPR$\downarrow$ & prec$_{1\%}$ & ms/img$\downarrow$ & fit s$\downarrow$ \\
    \midrule
    T0 all normal & 0.442 & 0.500 & 0 & 0 & 0 & -- & -- & -- \\
    T1 template & 0.702 & 0.762 & 0.693 & 0.767 & 0.465 & 0.018 & \textbf{0.12}$^\ast$ & \textbf{0.1} \\
    T2 PaDiM & 0.801\,$\pm$\,0.007 & 0.828\,$\pm$\,0.008 & 0.733\,$\pm$\,0.007 & 0.779\,$\pm$\,0.006
      & 0.332\,$\pm$\,0.012 & 0.024\,$\pm$\,0.001 & 7.38 & 2.8 \\
    T3 PatchCore & \textbf{0.885}\,$\pm$\,0.002 & \textbf{0.902}\,$\pm$\,0.002 & \textbf{0.806}\,$\pm$\,0.009
      & \textbf{0.808}\,$\pm$\,0.008 & \textbf{0.174}\,$\pm$\,0.021 & \textbf{0.046}\,$\pm$\,0.005 & 11.72 & 114 \\
    \bottomrule
  \end{tabular}
  \caption{Test results averaged over the 12 objects, mean\,$\pm$\,std over 3 seeds (T0 and T1 are
    deterministic); best per column in bold, T0 excluded. prec$_{1\%}$: Eq.~\ref{eq:prec-pi} at
    $\pi = 1\%$ from the object-averaged TPR and FPR of each seed (per object it is unreliable: with
    50--101 test normals, FPR is often 0). ms / img: median GPU latency incl.\ feature extraction ($^\ast$T1 on CPU, no backbone);
    fit: total over 12 objects. Without cashew, pcb1 and pcb2, PR-AUC is 0.655 / 0.785 / 0.874 for
    T1 / T2 / T3.}
  \label{tab:main}
\end{table*}

\paragraph{Main results.}
Tier~3 beats Tier~2 (Table~\ref{tab:main}) on 11 of 12 objects by $+0.084$ PR-AUC, a significant gap:
paired bootstrap 95\% CI $[0.065, 0.102]$ (test images resampled within object and label, 1,000
rounds), Wilcoxon $p = 0.003$ over objects, paired $t$-test $p = 0.003$ over seeds.
Without the three flagged objects it is $+0.089$ $[0.067, 0.109]$, $p = 0.012$.
Tier~2 beats Tier~1 by $+0.099$ but only on 9 of 12 objects ($p = 0.077$), a trend rather than a
significant gain.
Our Tier~3 ROC-AUC (90.2) is 2.2 points below the 92.4 reported on the official split
\citep{zou2022spot}, plausible for a re-implementation on a different test set.

\paragraph{Why Tier 2 loses.}
Tier~2 expects each part of the object at a fixed location, whereas Tier~3 matches patches from any
position; this robustness costs 41$\times$ the training time (114\,s vs.\ 2.8\,s, paid once) and
4.3\,ms more per image.
F2 hides the gap because the test set is 43\% anomalous (Tier~1 reaches F2 0.767 at FPR 0.465); at a
1\% prevalence only 2.4\% of Tier~2 and 4.6\% of Tier~3 alerts would be real defects.

\begin{table}[t]
  \centering
  \small
  \setlength{\tabcolsep}{3pt}
  \begin{tabular}{@{}llcc@{}}
    \toprule
    Factor & Setting & PR-AUC & $\Delta$ \\
    \midrule
    -- & full Tier 2 & 0.801\,$\pm$\,0.007 & -- \\
    \midrule
    Scoring & Euclidean 1-NN & 0.734\,$\pm$\,0.003 & $-0.067^\ast$ \\
            & cosine 1-NN    & 0.686\,$\pm$\,0.021 & $-0.115^\ast$ \\
    Contam. & 1\% of train & 0.797\,$\pm$\,0.007 & $-0.004$ \\
                  & 2\% of train & 0.794\,$\pm$\,0.008 & $-0.007$ \\
    Aggregation & top-1\% mean          & 0.799\,$\pm$\,0.006 & $-0.002$ \\
                & smoothed max &0.796\,$\pm$\,0.007 & $-0.004$ \\
    Features & global pooled & 0.771\,$\pm$\,0.006 & $-0.029^\ast$ \\
    \bottomrule
  \end{tabular}
  \caption{Tier~2 ablations, one factor at a time (12 objects, 3 seeds). $^\ast$: $|\Delta|$ exceeds
    the seed std of both settings. Full Tier~2: Mahalanobis, 0\% contamination, max, patch-level; smoothed max uses
    $\sigma = 4$.}
  \label{tab:ablation}
\end{table}

\paragraph{Controlled ablations.}
In Table~\ref{tab:ablation} the location-wise covariance matters most: a nearest-neighbour distance at
the same location costs 0.07--0.12 PR-AUC.
Pooling the map into one global vector costs 0.029 (0.038 on nine objects), in line with H2 (\S
ef{sec:analysis}).
The threshold rule leaves PR-AUC unchanged but moves the operating point beyond seed std: from best-F2
(recall / FPR / prec$_{1\%}$ = 0.82 / 0.33 / 0.024) to best-F1 (0.64 / 0.18 / 0.035), a top-10\% review
budget (0.35 / 0.04 / 0.082) or the 99th percentile of val normals (0.29 / 0.02 / 0.138).
Two results are negative.
Adding 1--2\% anomalies to train (4--16 images per object; our pool of 20 per object rules out 5\% and
10\%) changes PR-AUC by less than the seed std, likely because a median defect covers only 0.20\% of
the image, so far fewer than 2\% of \emph{locations} are contaminated.
Aggregation and the smoothing of \citet{defard2021padim} have no effect either, so H3 is not supported.
